\subsection{Transformation law and recurrence}
Use the standard global additive character, self-dual measures,
and the Gaussian $e^{-2\pi Q_B}$ at infinity.  Weil's operator
formulas \cite[I.13]{weil} give multiplication by $\mathbf e_p(bQ_B)$ for
$n(b)$, scaled pullback $f_p(x)\mapsto\gamma_{Q_B}(u)f_p(ux)$ for
$m(u)$ with $u$ a unit, and a constant times Fourier transform
for the Weyl element.  The upper unipotents $n(\mathbb Z_p)$ fix
$f_p$, because $Q_B$ is integral on its support.  If $L'_p$ is any
lattice of translations fixing $f_p$, its Fourier transform has
support in $(L'_p)^\vee$.  Choose $A_p$ with
$p^{A_p}Q_B((L'_p)^\vee)\subseteq\mathbb Z_p$.  Fourier conjugation
then proves that $n^-(p^{A_p}\mathbb Z_p)$ fixes $f_p$.
This includes $p=2$. At two take $A_2\ge2$ and use the comparison
with the ordinary theta multiplier given below; no splitting of
the metaplectic cover over all diagonal units is asserted.
For clarity, the local normalizations and the lower-unipotent
argument are as follows. Take
$\mathbf e_\infty(t)=e^{2\pi it}$ and
$\mathbf e_p(t)=e^{-2\pi i\{t\}_p}$, so that
$\prod_v\mathbf e_v(r)=1$ for $r\in\mathbb Q$.
For $Q(x)=x^{\mathsf T}Ax$ put
$B_Q(x,y)=Q(x+y)-Q(x)-Q(y)=2x^{\mathsf T}Ay$ and set
\[
 \mathcal F_Q\phi(y)=\int\phi(x)\mathbf e_p(B_Q(x,y))\,dx_Q,
 \qquad dx_Q=|\det(2A)|_p^{1/2}\,dx.
\]
Here $dx$ is the coordinate product measure with
$\operatorname{vol}(\mathbb Z_p^3)=1$.
At the real place use the same formula
$dx_Q=|\det(2A)|_\infty^{1/2}dx$ with coordinate Lebesgue measure.
Thus $\mathcal F_Q^2\phi(x)=\phi(-x)$.
If $\phi(x+\lambda)=\phi(x)$ for $\lambda\in L'$, translation
in this integral shows that its Fourier transform vanishes unless
$\mathbf e_p(B_Q(\lambda,y))=1$ for every $\lambda\in L'$.
This is precisely the asserted support in $(L')^\vee$.
Since this dual lattice is compact, an integer $A_p$ with
$p^{A_p}Q((L')^\vee)\subset\mathbb Z_p$ exists. The operator
for $n^-(c)$ is a Fourier conjugate of multiplication by
$\mathbf e_p(-cQ)$; consequently it fixes $\phi$ whenever
$c\in p^{A_p}\mathbb Z_p$. A scalar used to normalize the Weyl
operator cancels with its inverse in this conjugation.

All the specified local functions have such a translation lattice.
At an odd level prime the unit-root condition and its character
depend on a finite residue quotient. At two the square-unit
condition is a congruence modulo eight; the integral projectors
selecting the cyclic quotient, and the finite root character,
again depend on a finite residue quotient. At a prime dividing
$m$, valuation one and the nonzero-nilpotent sign depend on finite
residue data. The division-prime function is already a residue
function. Each support is contained in a lattice and is a finite
union of cosets of a sufficiently small lattice. This proves the
Schwartz--Bruhat and translation assertions for the actual weights.

For the reference $Q_0=\sum X_i^2$, the dyadic dual of
$\mathbb Z_2^3$ is $\tfrac12\mathbb Z_2^3$ and
$\operatorname{vol}_{dx_{Q_0}}(\mathbb Z_2^3)=2^{-3/2}$.
In particular its Fourier transform is
$2^{-3/2}1_{\frac12\mathbb Z_2^3}$, not $1_{\mathbb Z_2^3}$.
The lower-unipotent conclusion nevertheless holds at depth two,
because $4Q_0(\frac12\mathbb Z_2^3)\subset\mathbb Z_2$.
These positive measure factors affect neither the diagonal
eigenvalue quotient nor the counting coefficients of the theta sum.


